\chapter{Conclusiones y trabajo futuro}
\label{chapter:conclu}

\section{Conclusiones}

En este trabajo se desarrolló un pipeline de estimación de edad cerebral a partir de datos multimodales de neuroimagen, combinando conectividad funcional derivada de fMRI en reposo con medidas estructurales de materia gris obtenidas mediante imágenes T1w. El enfoque propuesto integra un $\beta$-VAE para la reducción de dimensionalidad no lineal de las matrices de conectividad funcional con un modelo XGBoost para la predicción supervisada de la edad, y fue evaluado sobre una cohorte de $N = 1245$ sujetos de la Red Latinoamericana de Investigación en Neurodegeneración (RedLaT).

El modelo principal alcanzó un MAE de 5,62 años y un $R^2$ de 0,411 sobre un conjunto de test independiente, valores que se encuentran dentro del rango reportado en la literatura para datasets de tamaño comparable. Cabe destacar que la mayoría de los estudios previos entrenan sus modelos exclusivamente sobre controles sanos, mientras que en este trabajo el pipeline fue entrenado sobre la cohorte completa, incluyendo pacientes con AD y FTD. A pesar de la mayor heterogeneidad que esto introduce, el modelo alcanza un rendimiento comparable, lo que sugiere que el pipeline es capaz de capturar patrones generales de envejecimiento cerebral aun en presencia de variabilidad patológica. Estos resultados confirman que es posible estimar la edad cerebral de forma razonable a partir de datos multimodales recolectados en un contexto multi-céntrico latinoamericano.

Uno de los hallazgos más relevantes del trabajo es que la integración de ambas modalidades mejora la predicción respecto al uso de cualquiera de ellas por separado. El estudio de ablaciones mostró que agregar los embeddings del VAE a los datos T1w reduce el MAE en 0,27 años, lo que indica que la conectividad funcional captura patrones de envejecimiento que los volúmenes de materia gris no contienen. Al mismo tiempo, los embeddings del VAE por sí solos resultan insuficientes como única fuente de información, lo que refuerza la importancia de la estrategia multimodal.

La comparación entre PCA y el $\beta$-VAE reveló un resultado notable: ambos métodos alcanzan un rendimiento prácticamente idéntico cuando se combinan con XGBoost optimizado, con diferencias menores a 0,02 en todas las métricas. Esto indica que, con un regresor suficientemente potente, las componentes principales capturan tanta información relevante para la predicción de edad como las representaciones no lineales del VAE. La elección del $\beta$-VAE siguió la tendencia predominante en la literatura de neuroimagen computacional, y la exploración de su espacio de hiperparámetros fue necesariamente limitada; no puede descartarse que configuraciones no evaluadas logren un rendimiento superior, particularmente con cohortes de mayor tamaño. Sin embargo, la reducción de dimensionalidad sí es necesaria: la comparación con los features crudos de FC demostró que el pipeline sin compresión obtiene un MAE de 6,00 frente a 5,62. Es decir, el paso clave es comprimir las 6670 dimensiones a 64, independientemente de si la compresión es lineal o no lineal.

Otro aspecto relevante del diseño fue la decisión de optimizar los hiperparámetros del VAE utilizando la predicción de edad como objetivo downstream, en lugar de minimizar exclusivamente el error de reconstrucción. Se espera que esta estrategia favorezca representaciones latentes orientadas a capturar información relevante al envejecimiento, en lugar de priorizar una reconstrucción fiel de las matrices FC que podría incluir variabilidad no relacionada con la edad. Verificar esta hipótesis mediante una comparación directa entre ambos criterios de optimización constituye una línea de trabajo futuro interesante.

El análisis del impacto de variables demográficas reveló que el sitio de adquisición es el factor con mayor influencia sobre la brain age gap, con diferencias de casi 8 años entre centros, lo que subraya la necesidad de técnicas de armonización en cohortes multicéntricas. Los años de educación mostraron una correlación significativa con la gap en pacientes con AD, posiblemente reflejando un efecto de reserva cognitiva por el cual los pacientes más educados acceden al diagnóstico en estadios más avanzados. En controles sanos, la tendencia fue negativa pero no significativa, consistente con un posible efecto neuroprotector de la educación. La inclusión de educación y sitio como features del modelo logró reducir el MAE a 5,17 años ---la mejor configuración evaluada en este trabajo---, lo que demuestra que estas variables contienen información predictiva complementaria a las features de neuroimagen. Estos resultados son particularmente relevantes en el contexto latinoamericano, donde la variabilidad en los niveles educativos y en las condiciones de adquisición es mayor que en cohortes más homogéneas. Cabe destacar que el pipeline es modular y reutilizable: su estructura facilita la incorporación de nuevas variables demográficas o la adaptación a distintas cohortes.

Finalmente, el experimento de entrenamiento exclusivo en controles sanos puso de manifiesto una limitación del tamaño de la muestra: con solo 473 sujetos CN, el modelo produjo señales de brain age gap débiles e inconsistentes entre los grupos clínicos, insuficientes para funcionar como biomarcador. Se realizó además una optimización Optuna dedicada para descartar que la limitación fuera atribuible a hiperparámetros subóptimos, confirmando que el problema es el tamaño muestral. No obstante, un resultado alentador surgió al incluir variables demográficas en el modelo CN-only: las gaps clínicas resultaron consistentemente positivas en ambos grupos (la dirección esperada en neurodegeneración), lo que sugiere que la combinación de variables demográficas con neuroimagen podría facilitar la construcción de un biomarcador cuando se disponga de cohortes CN más amplias. El pipeline diseñado es reutilizable para este fin: basta con filtrar los datos de entrenamiento a controles sanos y ejecutar el mismo flujo de optimización y evaluación.

\section{Trabajo futuro}

A partir de las limitaciones identificadas y los resultados obtenidos, se plantean varias líneas de trabajo que podrían extender y mejorar el enfoque propuesto.

La limitación más directa de este trabajo es el tamaño de la cohorte. Incrementar el número de sujetos, particularmente de controles sanos, permitiría no solo mejorar la generalización del modelo sino también habilitar el cálculo confiable de la brain age gap como biomarcador clínico. Los resultados del modelo CN-only con variables demográficas ---donde gaps consistentemente positivas en AD y FTD emergieron con solo 337 sujetos CN--- sugieren que, con cohortes de $N > 2000$ CN, la combinación de neuroimagen con datos demográficos podría producir un biomarcador robusto. En la misma dirección, la aplicación de técnicas de armonización entre sitios (como ComBat \cite{fortin2017harmonization}) resulta especialmente relevante: el análisis de variables demográficas demostró que el sitio de adquisición es el mayor confundidor de la gap, con diferencias de hasta 8 años entre centros. La armonización reduciría esta variabilidad espuria y podría mejorar tanto la precisión del modelo como la interpretabilidad de la brain age gap como indicador clínico. Además, la extensión del análisis demográfico a variables no disponibles o no incorporadas en este trabajo ---como nivel de ingresos, ocupación, estatus socioeconómico y edad de inicio de la enfermedad (ao)--- permitiría una caracterización más completa de los factores que modulan el envejecimiento cerebral en poblaciones latinoamericanas; en cohortes futuras, la recolección sistemática de estas variables facilitaría su inclusión tanto en el modelo predictivo como en el análisis de la brain age gap.

Desde el punto de vista del preprocesamiento, sería interesante evaluar el pipeline con atlas de mayor resolución que el AAL de 116 regiones. Parcelaciones de mayor resolución, como la de Schaefer et al.\ \cite{schaefer2018local} (disponible con hasta 1000 regiones), permitirían explorar si una mayor granularidad espacial captura patrones de conectividad asociados al envejecimiento que el AAL no detecta, aunque a costa de una mayor dimensionalidad que el VAE debería manejar.

En cuanto a la arquitectura del modelo, una extensión natural sería explorar enfoques de deep learning end-to-end que integren la representación y la predicción en una única red entrenada de forma conjunta. Las redes de grafos (GNN) aplicadas directamente sobre las matrices de conectividad, o las redes convolucionales 3D sobre los volúmenes T1w, podrían potencialmente mejorar el rendimiento, aunque típicamente requieren muestras mayores para entrenarse de forma efectiva.

Otra extensión directa consiste en ampliar las features estructurales incorporando medidas de volumen de materia blanca y de líquido cefalorraquídeo (LCR) por región, que también reflejan procesos de envejecimiento cerebral y podrían complementar los volúmenes de materia gris actualmente utilizados.

Por último, la incorporación de datos longitudinales (múltiples sesiones por sujeto) abriría la posibilidad de estudiar trayectorias individuales de envejecimiento, como han demostrado Elliott et al.\ \cite{elliott2019brain} al encontrar que la brain age gap en la edad adulta se asocia con el envejecimiento biológico acelerado y el deterioro cognitivo a lo largo del tiempo, mientras que la inclusión de otras modalidades de neuroimagen, como datos de difusión (DTI) o medidas de grosor cortical, podría aportar fuentes de información complementarias a las ya utilizadas.